\documentclass[12pt,a4paper,openany]{ctexbook}
% 基础宏包
\usepackage{geometry}
\usepackage{graphicx}
\usepackage{amsmath,amssymb}
\usepackage{booktabs}
\usepackage{array}
\usepackage{multirow}
\usepackage{enumitem}
\usepackage{xcolor}
\usepackage{listings}
\usepackage{hyperref}
\usepackage{fancyhdr}
\usepackage{tikz}
\usetikzlibrary{shapes,arrows,positioning,calc,decorations.pathreplacing}
% 页面设置
\geometry{left=2.5cm,right=2.5cm,top=2.5cm,bottom=2.5cm,headheight=15pt}
% 减少章节前后空白
\usepackage{titlesec}
\titlespacing*{\chapter}{0pt}{-20pt}{20pt}
\titlespacing*{\section}{0pt}{10pt}{5pt}
\titlespacing*{\subsection}{0pt}{8pt}{4pt}
% 超链接设置
\hypersetup{colorlinks=true,linkcolor=blue,urlcolor=blue}
% 代码样式
\lstset{basicstyle=\ttfamily\small,breaklines=true,frame=single,backgroundcolor=\color{gray!10}}
% 页眉页脚
\pagestyle{fancy}
\fancyhf{}
\fancyhead[L]{\leftmark}
\fancyhead[R]{\thepage}
\renewcommand{\headrulewidth}{0.4pt}
% 定义颜色盒子
\usepackage{tcolorbox}
\tcbuselibrary{skins,breakable}
\newtcolorbox{notebox}{colback=blue!5,colframe=blue!75!black,title=注意,fonttitle=\bfseries,breakable}
\newtcolorbox{tipbox}{colback=green!5,colframe=green!75!black,title=提示,fonttitle=\bfseries,breakable}
\newtcolorbox{formulabox}{colback=yellow!5,colframe=orange!75!black,title=公式,fonttitle=\bfseries,breakable}
\newtcolorbox{examplebox}[1][]{colback=gray!5,colframe=gray!75!black,title=#1,fonttitle=\bfseries,breakable}

\begin{document}
%========================================
\chapter{计算机网络填空题知识清单}
%========================================

\section{第一部分：网络基础概念（1-50题）}

1. 协议三要素为\textbf{语义}、\textbf{语法}和\textbf{时序}。

2. TCP/IP模型从低到高依次为\textbf{网络接口层}、\textbf{网间网层}、\textbf{传输层}、\textbf{应用层}。

3. 根据作用范围可以把网络分类为\textbf{LAN}、\textbf{MAN}、\textbf{WAN}、\textbf{PAN}。

4. 局域网的英文缩写为\textbf{LAN}，城域网为\textbf{MAN}，广域网为\textbf{WAN}。

5. OSI参考模型分为\textbf{7}层。

6. OSI模型从低到高依次是\textbf{物理层}、\textbf{数据链路层}、\textbf{网络层}、\textbf{传输层}、\textbf{会话层}、\textbf{表示层}、\textbf{应用层}。

7. 交换技术分为\textbf{电路交换}、\textbf{报文交换}、\textbf{分组交换}三种。

8. 计算机网络按带宽分为\textbf{窄带网}和\textbf{宽带网}。

9. 两个最重要的网络参考模型是\textbf{OSI/RM}模型和\textbf{TCP/IP}模型。

10. 计算机网络是现代\textbf{计算机}技术和\textbf{通信}技术密切结合的产物。

11. 在某层上进行通信所使用的规则的集合称为\textbf{该层的协议}。

12. 网络各层协议按层次顺序排列而成的协议称为网络的\textbf{协议栈}。

13. 对一个网络进行层次结构的划分时应做到：各层功能\textbf{明确}、相互独立。

14. 层间接口应\textbf{清晰}，穿越接口信息量尽可能少。

15. ISO的全称是\textbf{国际标准化组织}。

16. 计算机网络采用\textbf{分层}结构。

17. 每层完成一定的\textbf{功能}，每层都向它的上层提供一定的\textbf{服务}。

18. 低层协议对高层而言是\textbf{透明}的。

19. 相邻两层之间为\textbf{层间接口}。

20. 对等层用户通话时所必须遵守的规则称为\textbf{对等层协议}。

21. 会话层管理\textbf{不同主机上各进程间会话}。

22. 网络层传送的数据单元是\textbf{分组/数据报}。

23. 传输层传送的数据单元是\textbf{报文段}。

24. 数据链路层传送的数据单元是\textbf{帧}。

25. 物理层传送的数据单元是\textbf{比特}。

26. 分组交换重点解决的问题是\textbf{路由选择}、\textbf{流量控制}、\textbf{差错控制}。

27. 电路交换的特点是在通信之前需要建立一条\textbf{专用通道}。

28. 报文交换采用\textbf{存储转发}方式传输数据。

29. 分组交换将报文分成若干个\textbf{分组}进行传输。

30. 网络按使用者分为\textbf{公用网}和\textbf{专用网}。

31. Internet的应用层涵盖OSI/RM中的高\textbf{3}层。

32. 网络核心部分起特殊作用的是\textbf{路由器}。

33. 路由器是一种\textbf{专用}计算机。

34. 路由器的主要功能是实现\textbf{路由}和\textbf{转发}。

35. 在计算机网络中，带宽表示\textbf{单位时间内从网络中某一点到另一点所能通过的最高数据率}。

36. 带宽的单位是\textbf{bps}。

37. 时延包括\textbf{发送时延}、\textbf{传播时延}、\textbf{处理时延}、\textbf{排队时延}。

38. 发送时延等于\textbf{数据长度}除以\textbf{发送速率}。

39. 传播时延等于\textbf{传输距离}除以\textbf{传播速度}。

40. 往返时间RTT是指\textbf{从发送方发送数据到收到确认的时间}。

41. 信道利用率是指\textbf{信道被利用的时间占总时间的比例}。

42. 协议是\textbf{水平}的，服务是\textbf{垂直}的。

43. OSI模型中，N层为N+1层提供\textbf{服务}。

44. TCP/IP模型的网络接口层对应OSI的\textbf{物理层}和\textbf{数据链路层}。

45. TCP/IP模型的应用层对应OSI的\textbf{会话层}、\textbf{表示层}和\textbf{应用层}。

46. 虚电路服务需要在通信前建立\textbf{连接}。

47. 数据报服务中，每个分组独立选择\textbf{路由}。

48. 因特网采用的是\textbf{数据报}服务。

49. 表示层的功能包括数据\textbf{压缩}、\textbf{加密}和\textbf{格式转换}。

50. 计算机网络的定义为\textbf{一个互连的自主的计算机集合}。

\section{第二部分：物理层（51-100题）}

51. 物理层的主要任务所描述的4个特性为\textbf{机械特性}、\textbf{电气特性}、\textbf{功能特性}和\textbf{规程特性}。

52. 数据通信的几个主要指标有\textbf{信号传输速率}、\textbf{数据传输速率}、\textbf{信道容量}、\textbf{误码率}。

53. PCM的三个步骤是：\textbf{采样}、\textbf{量化}、\textbf{编码}。

54. 调制的基本方法有\textbf{幅移键控(ASK)}、\textbf{频移键控(FSK)}、\textbf{相移键控(PSK)}。

55. 传输介质分为\textbf{有线}和\textbf{无线}两大类。

56. 有线介质分为\textbf{同轴电缆}、\textbf{双绞线}、\textbf{光纤}。

57. 同轴电缆分为\textbf{基带同轴电缆}和\textbf{宽带同轴电缆}。

58. 无线传输介质主要有：\textbf{微波}、\textbf{红外线}、\textbf{激光}。

59. 网络拓扑结构有\textbf{星形}、\textbf{环形}、\textbf{树形}、\textbf{总线形}、\textbf{网形}。

60. 数字信号模拟传输时，数字信号变成音频信号的过程称为\textbf{调制}。

61. 音频信号变成数字信号的过程称为\textbf{解调}。

62. EIA RS-232C电气特性规定逻辑"1"的电平为\textbf{-5V}至\textbf{-15V}。

63. EIA RS-232C电气特性规定逻辑"0"的电平为\textbf{+5V}至\textbf{+15V}。

64. RS-232C采用\textbf{负}逻辑电平。

65. RS-232C中正负5V之间为\textbf{过渡区}，不作任何定义。

66. EIA RS-232C实际包括\textbf{主信道}和\textbf{辅信道}两信道。

67. 常用的多路复用技术包括：\textbf{FDM}、\textbf{TDM}、\textbf{WDM}、\textbf{CDMA}。

68. FDM的全称是\textbf{频分多路复用}。

69. TDM的全称是\textbf{时分多路复用}。

70. WDM的全称是\textbf{波分多路复用}。

71. CDMA的全称是\textbf{码分多址}。

72. 将一种数据形式转换成适合于在信道上传输的某种电信号形式，这类技术统称为\textbf{调制}。

73. 物理层的功能是实现原始数据在\textbf{物理媒体}上的传输。

74. 物理层为上一层提供传输原始\textbf{比特流}的物理连接。

75. DTE是指\textbf{数据终端设备}。

76. DCE是指\textbf{数据通信设备}。

77. DTE具有协议控制的\textbf{数据终端电路}的功能。

78. DCE为用户设备提供\textbf{入网}连接点。

79. EIA RS-232C机械特性规定了使用一个\textbf{25}芯标准连接器。

80. 香农公式为\textbf{$C=B\times\log_2(1+S/N)$}。

81. 奈氏准则公式为\textbf{$C=2B\log_2M$}。

82. 信噪比通常用\textbf{分贝(dB)}表示。

83. 信噪比的计算公式为\textbf{$10\log_{10}(S/N)$}。

84. 光纤是由能传导光波的\textbf{石英玻璃纤维}和\textbf{保护层}构成。

85. 光纤分为\textbf{单模光纤}和\textbf{多模光纤}两种。

86. 单模光纤比多模光纤传输距离\textbf{远}，传输速率\textbf{高}。

87. 双绞线传输介质是把两根导线绞在一起，这样可以减少\textbf{信号之间的相互串扰}。

88. 基带传输是指在信道中直接传输\textbf{数字}信号。

89. 频带传输是指将数字信号调制成\textbf{模拟}信号后传输。

90. 曼彻斯特编码中，比特1表示为码元中间从\textbf{高}到\textbf{低}跳变。

91. 曼彻斯特编码中，比特0表示为码元中间从\textbf{低}到\textbf{高}跳变。

92. 差分曼彻斯特编码中，码元中间始终有\textbf{跳变}。

93. 差分曼彻斯特编码中，比特0在码元开始处\textbf{有}跳变。

94. 差分曼彻斯特编码中，比特1在码元开始处\textbf{无}跳变。

95. 模拟信号数字化的过程包括\textbf{采样}、\textbf{量化}、\textbf{编码}。

96. 根据奈奎斯特采样定理，采样频率应大于等于信号最高频率的\textbf{2}倍。

97. 量化级数为128时，每个采样需要\textbf{7}位编码。

98. 量化级数为256时，每个采样需要\textbf{8}位编码。

99. 电话信道的带宽约为\textbf{3400}Hz。

100. 中继器工作在OSI模型的\textbf{物理}层。

\section{第三部分：数据链路层（101-150题）}

101. 流量控制是指\textbf{使得发送端的发送速率不要太快，接收端来得及接收}。

102. 允许发送站发送多个帧而不需要马上应答，这就是\textbf{滑动窗口}协议。

103. HDLC有三种不同类型的帧，分别称为\textbf{信息帧}、\textbf{监督帧}、\textbf{无编号帧}。

104. PPP协议包括三个部分：\textbf{PPP}、\textbf{NCP}和\textbf{LCP}。

105. 使用\textbf{RARP}将物理地址转换为IP地址。

106. 使用\textbf{ARP}将IP地址转换为物理地址。

107. 数据链路层的协议数据单元为\textbf{帧}。

108. 以太网采用的介质访问控制协议是\textbf{CSMA/CD}。

109. CSMA/CD的工作原理可以概括为：\textbf{先听后说}、\textbf{边说边听}、\textbf{冲突停止}、\textbf{随机重发}。

110. 以太网采用的退避算法是\textbf{二进制指数截断退避算法}。

111. MAC地址长度为\textbf{48}位，即\textbf{6}字节。

112. 10Mbps以太网的争用期为\textbf{51.2$\mu$s}。

113. 10Mbps以太网的最小帧长为\textbf{64}字节。

114. 交换机工作在\textbf{数据链路}层。

115. 路由器工作在\textbf{网络}层。

116. 集线器工作在\textbf{物理}层。

117. 交换机的转发规则包括：\textbf{自学习}、\textbf{明确转发}、\textbf{泛洪}、\textbf{过滤}。

118. VLAN的作用是将物理局域网在逻辑上划分为多个\textbf{广播域}。

119. 802.1Q帧中VLAN标记的VID字段为\textbf{12}位。

120. VLAN ID的范围为\textbf{1-4094}。

121. 交换机端口类型主要有\textbf{Access}端口和\textbf{Trunk}端口。

122. 生成树协议STP用于解决冗余链路形成的\textbf{环路}问题。

123. 使用双绞线的10M以太网标准为\textbf{10Base-T}。

124. 使用双绞线的100M以太网标准为\textbf{100Base-T}。

125. 在Go-Back-N协议中，当序号为n比特时，发送窗口最大为\textbf{$2^n-1$}。

126. 在选择重传协议中，当序号为n比特时，发送窗口最大为\textbf{$2^{n-1}$}。

127. 停止等待协议的发送窗口大小为\textbf{1}。

128. 停止等待协议的接收窗口大小为\textbf{1}。

129. Go-Back-N协议的接收窗口大小为\textbf{1}。

130. CRC校验是一种\textbf{差错}检测方法。

131. CRC校验基于\textbf{多项式}除法原理。

132. PPP帧的标志字段值为\textbf{7E}。

133. PPP协议使用\textbf{字节}填充实现透明传输。

134. HDLC是面向\textbf{比特}的通信协议。

135. PPP协议的工作状态包括\textbf{链路静止}、\textbf{链路建立}、\textbf{鉴别}、\textbf{网络层协议}。

136. 以太网的广播地址为\textbf{FF-FF-FF-FF-FF-FF}。

137. MAC地址的前3字节表示\textbf{OUI（厂商标识）}。

138. MAC地址的后3字节表示\textbf{设备标识}。

139. 局域网中访问冲突的根源是\textbf{共享}介质。

140. 以太网帧的最大长度为\textbf{1518}字节。

141. 以太网帧的最小长度为\textbf{64}字节。

142. 以太网帧的数据字段最大为\textbf{1500}字节。

143. 以太网帧的数据字段最小为\textbf{46}字节。

144. 比特填充是在连续\textbf{5}个1后插入一个0。

145. 争用期也称为\textbf{碰撞窗口}。

146. 争用期等于\textbf{2}倍的单向传播时延。

147. 最小帧长等于争用期乘以\textbf{数据率}。

148. 802.1Q帧的TPID字段值为\textbf{0x8100}。

149. STP中优先级最高的交换机称为\textbf{根桥}。

150. 非根桥上到根桥路径最短的端口称为\textbf{根端口}。

\section{第四部分：网络层-IP地址（151-200题）}

151. IPv4地址长度为\textbf{32}位。

152. IPv6地址长度为\textbf{128}位。

153. A类IP地址的网络号占\textbf{8}位（实际可用7位）。

154. A类IP地址的主机号占\textbf{24}位。

155. B类IP地址的网络号占\textbf{16}位。

156. B类IP地址的主机号占\textbf{16}位。

157. C类IP地址的网络号占\textbf{24}位。

158. C类IP地址的主机号占\textbf{8}位。

159. B类IP地址网络允许容纳的最大主机数为\textbf{65534}。

160. C类IP地址网络允许容纳的最大主机数为\textbf{254}。

161. A类IP地址的默认子网掩码为\textbf{255.0.0.0}。

162. B类IP地址的默认子网掩码为\textbf{255.255.0.0}。

163. C类IP地址的默认子网掩码为\textbf{255.255.255.0}。

164. 私有IP地址中，A类为\textbf{10.0.0.0/8}。

165. 私有IP地址中，B类为\textbf{172.16.0.0/12}。

166. 私有IP地址中，C类为\textbf{192.168.0.0/16}。

167. IP地址\textbf{0.0.0.0}表示本网络本主机。

168. IP地址\textbf{255.255.255.255}表示本网段广播。

169. \textbf{127.x.x.x}地址段用于环回测试。

170. A类IP地址的首位为\textbf{0}。

171. B类IP地址的前两位为\textbf{10}。

172. C类IP地址的前三位为\textbf{110}。

173. D类IP地址的前四位为\textbf{1110}。

174. D类IP地址用于\textbf{多播}。

175. E类IP地址用于\textbf{保留}。

176. A类IP地址的范围是\textbf{1.0.0.0}到\textbf{126.x.x.x}。

177. B类IP地址的范围是\textbf{128.0.0.0}到\textbf{191.255.x.x}。

178. C类IP地址的范围是\textbf{192.0.0.0}到\textbf{223.255.255.x}。

179. 子网掩码的作用是区分IP地址中的\textbf{网络}部分和\textbf{主机}部分。

180. IP地址AND子网掩码等于\textbf{网络地址}。

181. CIDR的斜线记法表示为\textbf{IP地址/前缀长度}。

182. CIDR消除了传统的\textbf{A}、\textbf{B}、\textbf{C}类地址的界限。

183. 路由器转发时，选择与目的地址匹配的\textbf{最长前缀}的路由。

184. VLSM允许在同一网络中使用不同长度的\textbf{子网掩码}。

185. 主机号全0的IP地址表示\textbf{网络地址}。

186. 主机号全1的IP地址表示\textbf{直接广播地址}。

187. IPv4向IPv6进行过渡使用的两种策略是\textbf{双协议栈}和\textbf{隧道技术}。

188. IPv6地址使用\textbf{冒号}十六进制表示。

189. IPv6地址中连续的0可以用\textbf{::}表示。

190. NAT的全称是\textbf{网络地址转换}。

191. NAPT使用IP地址加\textbf{端口号}进行映射。

192. 为了解决IP地址耗尽的问题，治本的措施是\textbf{采用IPv6}。

193. 划分子网需要从\textbf{主机}号中借位。

194. 将192.168.1.0/24划分为4个子网，每个子网的前缀长度为\textbf{/26}。

195. /26的子网掩码为\textbf{255.255.255.192}。

196. /27的子网掩码为\textbf{255.255.255.224}。

197. /28的子网掩码为\textbf{255.255.255.240}。

198. /30的子网掩码为\textbf{255.255.255.252}。

199. /30子网可分配给主机的地址数为\textbf{2}个。

200. 路由聚合也称为\textbf{超网/构造超网}。

\section{第五部分：网络层-路由与协议（201-250题）}

201. 路由器数据层面的交换结构类型有\textbf{经内存交换}、\textbf{经总线交换}、\textbf{经互联网络交换}。

202. IP数据报首部中，TTL字段表示\textbf{生存时间}。

203. TTL每经过一个路由器减\textbf{1}。

204. IP数据报分片时，片偏移字段以\textbf{8}字节为单位。

205. 以太网的MTU为\textbf{1500}字节。

206. RIP协议以\textbf{跳数}作为度量标准。

207. RIP协议的最大跳数为\textbf{15}。

208. OSPF协议以\textbf{链路代价}作为度量标准。

209. OSPF协议使用\textbf{Dijkstra最短路径}算法。

210. RIP协议封装在\textbf{UDP}中。

211. OSPF协议直接封装在\textbf{IP}中。

212. BGP协议封装在\textbf{TCP}中。

213. \textbf{BGP}协议用于自治系统之间的路由选择。

214. RIP协议使用的端口号为\textbf{520}。

215. BGP协议使用的端口号为\textbf{179}。

216. ICMP协议用于在主机和路由器之间传递\textbf{控制信息}和\textbf{错误报告}。

217. ping命令使用ICMP的\textbf{回送请求/回答}报文。

218. traceroute利用\textbf{TTL}递增和ICMP时间超过报文发现路径。

219. IP数据报首部最小为\textbf{20}字节。

220. IP数据报首部最大为\textbf{60}字节。

221. IP首部中协议字段，TCP的值为\textbf{6}。

222. IP首部中协议字段，UDP的值为\textbf{17}。

223. IP首部中协议字段，ICMP的值为\textbf{1}。

224. IP首部校验和只校验\textbf{首部}。

225. IP数据报的最大长度为\textbf{65535}字节。

226. 分片时MF=1表示\textbf{后面还有分片}。

227. 分片时MF=0表示\textbf{最后一个分片}。

228. DF=1表示\textbf{不允许分片}。

229. 能够抑制广播风暴的网络设备是\textbf{路由器}。

230. RIP协议每\textbf{30}秒向相邻路由器广播路由表。

231. RIP协议中跳数=16表示\textbf{不可达}。

232. OSPF协议属于\textbf{链路状态}路由协议。

233. RIP协议属于\textbf{距离向量}路由协议。

234. 内部网关协议包括\textbf{RIP}和\textbf{OSPF}。

235. 外部网关协议主要是\textbf{BGP}。

236. ICMP差错报告报文包括\textbf{终点不可达}、\textbf{源点抑制}、\textbf{时间超过}、\textbf{参数问题}、\textbf{改变路由}。

237. VPN的全称是\textbf{虚拟专用网}。

238. VPN利用\textbf{IP隧道}技术封装内部数据报。

239. 直接交付是指目的主机与源主机在\textbf{同一}网络。

240. 间接交付需要将数据报发送给\textbf{网关/路由器}。

241. 默认路由的目的网络地址为\textbf{0.0.0.0}。

242. 默认路由的子网掩码为\textbf{0.0.0.0}。

243. IPv6首部长度固定为\textbf{40}字节。

244. IPv6取消了\textbf{首部校验和}字段。

245. IPv6不支持\textbf{广播}，用多播替代。

246. IPv6中只有\textbf{源主机}可以进行分片。

247. ARP请求报文使用\textbf{广播}方式发送。

248. ARP应答报文使用\textbf{单播}方式发送。

249. 代理ARP用于\textbf{不同}网络之间的通信。

250. 免费ARP用于检测\textbf{IP地址}冲突。

\section{第六部分：传输层（251-300题）}

251. 传输层提供\textbf{端到端}通信服务。

252. 端口号的范围是\textbf{0-65535}。

253. 熟知端口的范围是\textbf{0-1023}。

254. 登记端口的范围是\textbf{1024-49151}。

255. 动态端口的范围是\textbf{49152-65535}。

256. TCP协议是\textbf{面向连接}的传输协议（连接方式）。

257. TCP协议是\textbf{可靠}的传输协议（可靠性）。

258. UDP协议是\textbf{无连接}的传输协议（连接方式）。

259. UDP协议是\textbf{不可靠}的传输协议（可靠性）。

260. TCP首部最小\textbf{20}字节。

261. TCP首部最大\textbf{60}字节。

262. UDP首部固定\textbf{8}字节。

263. TCP建立连接需要\textbf{3}次握手。

264. TCP释放连接需要\textbf{4}次挥手。

265. TCP连接释放时，主动关闭方发送最后一个ACK后进入\textbf{TIME\_WAIT}状态。

266. TIME\_WAIT状态需要等待\textbf{2MSL}时间。

267. MSL的全称是\textbf{最大报文段生存时间}。

268. TCP的拥塞控制算法包括\textbf{慢开始}、\textbf{拥塞避免}、\textbf{快重传}、\textbf{快恢复}。

269. TCP发送窗口等于min(\textbf{接收窗口rwnd}, \textbf{拥塞窗口cwnd})。

270. 收到\textbf{3}个重复ACK时，TCP执行快重传。

271. TCP序号字段表示\textbf{本报文段数据的第一个字节的序号}。

272. TCP确认号字段表示\textbf{期望收到的下一个字节的序号}。

273. TCP首部中SYN标志用于\textbf{建立连接}。

274. TCP首部中FIN标志用于\textbf{释放连接}。

275. TCP首部中ACK标志表示\textbf{确认号}有效。

276. TCP首部中RST标志用于\textbf{复位连接}。

277. TCP首部中PSH标志表示\textbf{尽快交付应用层}。

278. TCP首部中URG标志表示\textbf{紧急指针}有效。

279. TCP采用\textbf{累积}确认机制。

280. UDP支持\textbf{多播/广播}通信。

281. TCP是\textbf{点对点}通信。

282. 慢开始时cwnd从\textbf{1}开始。

283. 慢开始阶段cwnd呈\textbf{指数}增长。

284. 拥塞避免阶段cwnd呈\textbf{线性}增长。

285. 发生超时后ssthresh变为原cwnd的\textbf{一半}。

286. 发生超时后cwnd变为\textbf{1}。

287. 快恢复后cwnd变为\textbf{ssthresh}。

288. 零窗口探测报文用于解决\textbf{死锁}问题。

289. TCP使用\textbf{滑动窗口}机制进行流量控制。

290. 接收窗口大小由\textbf{接收方}通告。

291. 拥塞窗口大小由\textbf{发送方}维护。

292. 超时重传时间RTO根据\textbf{RTT}动态调整。

293. 三次握手的第一次，客户端发送\textbf{SYN}报文。

294. 三次握手的第二次，服务器发送\textbf{SYN+ACK}报文。

295. 三次握手的第三次，客户端发送\textbf{ACK}报文。

296. 四次挥手的第一次，主动方发送\textbf{FIN}报文。

297. 客户端发送SYN后进入\textbf{SYN\_SENT}状态。

298. 服务器收到SYN后进入\textbf{SYN\_RCVD}状态。

299. 连接建立后双方进入\textbf{ESTABLISHED}状态。

300. 被动关闭方收到FIN后进入\textbf{CLOSE\_WAIT}状态。

\section{第七部分：应用层协议（301-350题）}

301. HTTP的中文含义是\textbf{超文本传输协议}。

302. TCP的中文含义是\textbf{传输控制协议}。

303. UDP的中文含义是\textbf{用户数据报协议}。

304. FTP的中文含义是\textbf{文件传输协议}。

305. DNS的中文含义是\textbf{域名系统}。

306. SMTP的中文含义是\textbf{简单邮件传输协议}。

307. DHCP的中文含义是\textbf{动态主机配置协议}。

308. FTP使用\textbf{2}个TCP连接。

309. FTP控制连接端口为\textbf{21}。

310. FTP数据连接端口为\textbf{20}。

311. DNS主要使用\textbf{UDP}协议。

312. DNS的端口号为\textbf{53}。

313. HTTP默认端口为\textbf{80}。

314. HTTPS默认端口为\textbf{443}。

315. SMTP端口为\textbf{25}。

316. POP3端口为\textbf{110}。

317. IMAP端口为\textbf{143}。

318. SSH端口为\textbf{22}。

319. Telnet端口为\textbf{23}。

320. DHCP客户端使用端口\textbf{68}。

321. DHCP服务器使用端口\textbf{67}。

322. DHCP工作过程包括\textbf{Discover}、\textbf{Offer}、\textbf{Request}、\textbf{ACK}四个步骤。

323. \textbf{ipconfig}命令可以用来查看TCP/IP配置信息。

324. \textbf{ping}命令测试主机是否可达。

325. HTTP状态码200表示\textbf{成功（OK）}。

326. HTTP状态码404表示\textbf{未找到}。

327. HTTP状态码500表示\textbf{服务器内部错误}。

328. HTTP状态码301表示\textbf{永久移动}。

329. HTTP状态码302表示\textbf{临时移动}。

330. HTTP状态码403表示\textbf{禁止访问}。

331. HTTP是\textbf{无状态}协议。

332. HTTP基于\textbf{TCP}协议。

333. HTTP/1.0使用\textbf{非持久}连接。

334. HTTP/1.1默认使用\textbf{持久}连接。

335. HTTP请求方法GET用于\textbf{获取资源}。

336. HTTP请求方法POST用于\textbf{提交数据}。

337. HTTP请求方法PUT用于\textbf{上传文件}。

338. HTTP请求方法DELETE用于\textbf{删除资源}。

339. Cookie存储在\textbf{客户}端。

340. Session存储在\textbf{服务器}端。

341. DNS服务器类型包括\textbf{根域名服务器}、\textbf{顶级域名服务器}、\textbf{权威域名服务器}、\textbf{本地域名服务器}。

342. DNS查询方式有\textbf{递归查询}和\textbf{迭代查询}两种。

343. 本地DNS服务器与客户端之间通常采用\textbf{递归}查询。

344. 本地DNS服务器向根DNS服务器通常采用\textbf{迭代}查询。

345. 域名结构从右到左依次为\textbf{根域}、\textbf{顶级域}、\textbf{二级域}。

346. SMTP是\textbf{推送/发送}协议。

347. POP3和IMAP是\textbf{拉取/接收}协议。

348. 邮件从发送方服务器到接收方服务器使用\textbf{SMTP}协议。

349. FTP的主动模式由\textbf{服务器}主动连接数据端口。

350. FTP的被动模式由\textbf{客户端}主动连接数据端口。

\section{第八部分：应用层补充（351-400题）}

351. Web缓存可以减少\textbf{客户端请求}的响应时间。

352. Web缓存可以减少机构接入链路的\textbf{流量}。

353. 条件GET请求头包含\textbf{If-Modified-Since}字段。

354. 若对象未修改，服务器返回\textbf{304 Not Modified}状态码。

355. DNS缓存记录有\textbf{TTL}，到期后会被删除。

356. 顶级域名包括\textbf{通用}顶级域和\textbf{国家/地区}顶级域。

357. .com属于\textbf{通用}顶级域。

358. .cn属于\textbf{国家/地区}顶级域。

359. URL的组成包括\textbf{协议}、\textbf{主机名}、\textbf{路径}。

360. HTTPS使用\textbf{SSL/TLS}协议进行加密。

361. 电子邮件系统由\textbf{用户代理}、\textbf{邮件服务器}、\textbf{邮件协议}组成。

362. 用户代理是指\textbf{邮件客户端}软件。

363. MIME用于扩展电子邮件支持\textbf{多媒体}内容。

364. P2P模式中每个节点既是\textbf{客户端}又是\textbf{服务器}。

365. BitTorrent是一种\textbf{P2P}文件分发协议。

366. CDN的全称是\textbf{内容分发网络}。

367. CDN用于加速\textbf{内容}的分发。

368. SNMP用于\textbf{网络}管理。

369. SNMP使用\textbf{UDP}协议。

370. SNMP的端口号为\textbf{161}。

371. NTP用于\textbf{时间}同步。

372. NTP使用\textbf{UDP}协议。

373. NTP的端口号为\textbf{123}。

374. TFTP的全称是\textbf{简单文件传输协议}。

375. TFTP使用\textbf{UDP}协议。

376. TFTP的端口号为\textbf{69}。

377. RDP的默认端口为\textbf{3389}。

378. MySQL的默认端口为\textbf{3306}。

379. Redis的默认端口为\textbf{6379}。

380. LDAP的默认端口为\textbf{389}。

381. IMAP相比POP3的优点是邮件保留在\textbf{服务器}。

382. WebSocket用于实现\textbf{全双工}通信。

383. SSL/TLS工作在\textbf{传输}层和\textbf{应用}层之间。

384. 数字证书由\textbf{CA}颁发。

385. DHCP Discover报文的源IP为\textbf{0.0.0.0}。

386. DHCP Discover报文的目的IP为\textbf{255.255.255.255}。

387. DHCP分配的IP地址有\textbf{租约时间}，到期需要续租。

388. DNS的A记录用于将域名映射到\textbf{IPv4地址}。

389. DNS的AAAA记录用于将域名映射到\textbf{IPv6地址}。

390. DNS的CNAME记录用于设置\textbf{别名}。

391. DNS的MX记录用于指定\textbf{邮件}服务器。

392. DNS的NS记录用于指定\textbf{域名}服务器。

393. DNS的PTR记录用于\textbf{反向}解析。

394. DNS的TXT记录用于存储\textbf{文本}信息。

395. HTTP/2支持\textbf{多路}复用。

396. HTTP/2使用\textbf{HPACK}压缩首部。

397. HTTP/3基于\textbf{QUIC}协议。

398. QUIC协议基于\textbf{UDP}。

399. REST是一种\textbf{Web服务}架构风格。

400. JSON的全称是\textbf{JavaScript对象表示法}。

\section{第九部分：综合布线与网络设备（401-450题）}

401. 综合布线系统包括\textbf{工作区子系统}、\textbf{水平子系统}、\textbf{管理子系统}、\textbf{垂直干线子系统}、\textbf{设备间子系统}、\textbf{建筑群子系统}六个子系统。

402. 工作区子系统连接\textbf{终端设备}和\textbf{信息插座}。

403. 水平子系统连接\textbf{信息插座}和\textbf{配线架}。

404. 垂直干线子系统连接\textbf{各楼层配线架}和\textbf{主配线架}。

405. 设备间子系统放置\textbf{核心网络}设备。

406. 建筑群子系统连接\textbf{各建筑物}之间的通信。

407. 管理子系统用于\textbf{跳线管理}和\textbf{线缆测试}。

408. 水平布线一般使用\textbf{双绞线}电缆。

409. 垂直干线一般使用\textbf{光纤或大对数双绞线}电缆。

410. 建筑群之间一般使用\textbf{光纤}电缆。

411. 网络互连设备按层次分为\textbf{中继器}、\textbf{网桥/交换机}、\textbf{路由器}、\textbf{网关}。

412. 中继器用于\textbf{放大/再生}信号。

413. 集线器是多端口的\textbf{中继器}。

414. 网桥用于连接两个\textbf{局域网}。

415. 交换机是多端口的\textbf{网桥}。

416. 路由器用于连接不同的\textbf{网络}。

417. 网关用于连接不同\textbf{协议}的网络。

418. 防火墙用于\textbf{保护}网络安全。

419. 负载均衡器用于\textbf{分发}流量。

420. 三层交换机具有\textbf{路由}功能。

421. 集线器的所有端口共享\textbf{带宽/冲突域}。

422. 交换机的每个端口是独立的\textbf{冲突域}。

423. 集线器工作在\textbf{半双工}模式。

424. 交换机可以工作在\textbf{全双工}模式。

425. 路由器隔离\textbf{广播}域。

426. 交换机隔离\textbf{冲突}域。

427. 集线器不隔离\textbf{冲突}域。

428. 冲突域是指\textbf{可能发生冲突}的范围。

429. 广播域是指\textbf{广播帧能到达}的范围。

430. VLAN可以隔离\textbf{广播}域。

431. 交换机的转发方式有\textbf{直通转发}、\textbf{存储转发}、\textbf{无碎片转发}。

432. 直通转发延迟\textbf{小}，但不检查错误。

433. 存储转发延迟\textbf{大}，但可以检查错误。

434. 无碎片转发只检查帧的前\textbf{64}字节。

435. 路由器的主要功能是\textbf{路由选择}和\textbf{分组转发}。

436. 路由表包含\textbf{目的网络}、\textbf{下一跳}、\textbf{接口}等信息。

437. 静态路由由\textbf{管理员}手动配置。

438. 动态路由由\textbf{路由协议}自动学习。

439. 直连路由是路由器\textbf{直接}连接的网络。

440. 默认路由用于转发\textbf{无匹配路由}的数据包。

441. PoE的全称是\textbf{以太网供电}。

442. PoE可以通过\textbf{网线}供电。

443. 堆叠技术可以将多台交换机虚拟成\textbf{一}台。

444. 链路聚合可以将多条链路捆绑成\textbf{一}条。

445. LACP的全称是\textbf{链路聚合控制协议}。

446. 端口镜像用于\textbf{监控/抓取}流量。

447. ACL的全称是\textbf{访问控制列表}。

448. ACL用于\textbf{过滤}数据包。

449. QoS的全称是\textbf{服务质量}。

450. QoS用于保证\textbf{网络服务质量}。

\section{第十部分：公式与综合（451-500题）}

451. 香农公式为\textbf{$C=B\times\log_2(1+S/N)$}。

452. 奈氏准则公式为\textbf{$C=2B\log_2M$}。

453. 信噪比的计算公式为\textbf{$10\log_{10}(S/N)$ dB}。

454. 发送时延的计算公式为\textbf{数据长度/发送速率}。

455. 传播时延的计算公式为\textbf{传输距离/传播速度}。

456. 信道利用率的计算公式为\textbf{有效传输时间/总时间}。

457. 停等协议的信道利用率公式为\textbf{$T_D/(T_D+RTT+T_A)$}。

458. 以太网最小帧长的计算公式为\textbf{争用期$\times$数据率}。

459. IP分片的片偏移计算公式为\textbf{数据偏移量/8}。

460. 子网可分配主机数的计算公式为\textbf{$2^n-2$}。

461. 划分n个子网需要借\textbf{$\lceil\log_2 n\rceil$}位。

462. /24子网可分配\textbf{254}个主机地址。

463. /25子网可分配\textbf{126}个主机地址。

464. /26子网可分配\textbf{62}个主机地址。

465. /27子网可分配\textbf{30}个主机地址。

466. /28子网可分配\textbf{14}个主机地址。

467. /29子网可分配\textbf{6}个主机地址。

468. /30子网可分配\textbf{2}个主机地址。

469. 1Mbps等于\textbf{$10^6$}bps。

470. 1Gbps等于\textbf{1000}Mbps。

471. 1KB等于\textbf{1024}字节。

472. 1MB等于\textbf{1024}KB。

473. 1字节等于\textbf{8}比特。

474. 以太网帧间隙为\textbf{12}字节。

475. 以太网前导码为\textbf{8}字节。

476. TCP序号字段占\textbf{4}字节。

477. TCP确认号字段占\textbf{4}字节。

478. TCP窗口字段占\textbf{2}字节。

479. IP首部的版本字段占\textbf{4}位。

480. IP首部的首部长度字段占\textbf{4}位。

481. IP首部的总长度字段占\textbf{16}位。

482. IP首部的TTL字段占\textbf{8}位。

483. IP首部的协议字段占\textbf{8}位。

484. 以太网帧的目的MAC字段占\textbf{6}字节。

485. 以太网帧的源MAC字段占\textbf{6}字节。

486. 以太网帧的类型字段占\textbf{2}字节。

487. 以太网帧的FCS字段占\textbf{4}字节。

488. UDP首部的长度字段占\textbf{2}字节。

489. UDP首部的校验和字段占\textbf{2}字节。

490. CRC-32的校验码长度为\textbf{32}位。

491. 在CRC中，生成多项式$x^3+x+1$对应二进制为\textbf{1011}。

492. 在CRC中，生成多项式$x^4+x+1$对应二进制为\textbf{10011}。

493. 在CRC中，生成多项式$x^5+x^3+x^2+1$对应二进制为\textbf{101101}。

494. 100Base-TX使用\textbf{5}类双绞线。

495. 1000Base-T使用\textbf{5e/6}类双绞线。

496. 10GBase-T使用\textbf{6a/7}类双绞线。

497. 光纤的传输速率可达\textbf{Tbps}级别。

498. 5G网络的理论峰值速率可达\textbf{20Gbps}。

499. WiFi 6的标准名称是\textbf{802.11ax}。

500. 蓝牙5.0的理论传输速率为\textbf{2Mbps}。

\end{document}
